This research project examined lexical semantic change in ADHD and autism across thirteen years of general web discourse, motivated by the concern that the popularisation of mental health language erodes the meaning of the terms on which diagnosed people rely. It constructed a reproducible Common Crawl pipeline spanning 2014--2026, separated clinical from lived-experience framing with a validated hierarchical classifier, and estimated salience, sentiment, intensity, breadth, and thematic-neighbour trajectories using a refined adaptation of the SIBling framework, benchmarked against the framework's original operationalisations.

The hypothesised concept creep did not materialise. ADHD and autism did not drift into milder emotional contexts, and their contexts did not broaden; ADHD contexts contracted markedly, and the thematic core of both discourses---diagnosis, childhood, family---remained stable throughout. What changed was the composition of the discourse: a robust shift from clinical toward lived-experience framing, strongest for ADHD, together with a recent, tentative positification of clinical contexts and growing entanglement between the two concepts. In sampled open-web discourse, autism became less prominent rather than more, underscoring that the contemporary surge in ADHD and autism talk is unevenly distributed across the digital landscape.

The study makes four contributions. Substantively, it extends concept creep and therapy-speak research to ADHD and autism and provides evidence that two institutionally codified diagnostic concepts resisted the semantic loosening documented for ordinary-language harm concepts, with category expansion proceeding through application rather than through the dilution of meaning in use. Evidentially, it moves LSC research on mental-health concepts from curated corpora to general web discourse and shows that conclusions do not simply transfer between the two. Analytically, it demonstrates at web scale that discourse composition and semantic change must be separated, since frame-specific trajectories repeatedly diverged from aggregates. Methodologically, its robustness checks show that key SIBling conclusions are conditional on lexicon and encoder choice, identifying measurement validation as a priority for the field.

The broader implication is a reframing of the trivialisation concern. For ADHD and autism on the open web, the words have largely held their shape; it is the conversation around them---who speaks, in what frame, and on which platforms---that has been transformed. Whether the same holds where that conversation is now loudest remains the question this project leaves open, and the tools it provides are designed to answer it.
